\chapter[Normalization and Ordinal Analysis]{Normalization, Proof Terms, and Ordinal Analysis}
\label{ch:normalization-curry-howard}

\section{Natural-deduction normalization}
\label{sec:normalization-curry-howard-01}

The normalization theory of natural deduction developed in this chapter
originates with Prawitz \cite{Prawitz1965}.

\begin{definition}[Detour]
A detour is a segment of a natural-deduction derivation in which a formula is
introduced by an introduction rule and is immediately used as the major
premise of the matching elimination rule; the \emph{major premise} of an
elimination rule is the premise exhibiting the connective or quantifier
that the rule eliminates.  A derivation is normal when it has
no detour and satisfies the required permutation conditions for disjunction
and existential elimination.
\end{definition}

For implication, the principal reduction has the schematic form
\[
  \dfrac{
    \dfrac{\begin{matrix}{[A]}^i\\ \vdots\,D\\ B\end{matrix}}
          {A\to B}\;\to I_i
    \qquad
    \begin{matrix}\vdots\,E\\ A\end{matrix}
  }{B}\;\to E
  \quad\leadsto\quad
  \begin{matrix}\vdots\,E\\ A\\ \vdots\,D\\ B\end{matrix}\;,
\]
where \(D[E/A^i]\) substitutes the derivation \(E\) for every discharged
occurrence carrying label \(i\).  This is substitution of derivations, not term
substitution inside formulas.

For the theorem below, the calculus is exactly the finite-tree
natural-deduction presentation of \cref{ch:prop-deduction,ch:fol-deduction}
with the rules
\(\land I,\land E_1,\land E_2,\lor I_1,\lor I_2,\lor E,\to I,\to E,
\bottom E,\forall I,\forall E,\exists I,\exists E\).  It omits
\(\mathsf{RAA}\), the primitive-negation rules, and the equality rules, and
uses \(\neg A\defeq A\to\bottom\).  Its quantifier substitutions are the total
capture-avoiding substitutions fixed earlier, and its discharge and
eigenvariable conditions are literally those attached to the cited rules.
This declaration fixes the derivation objects to which ``reduces'' and
``normal'' refer; it does not posit an unnamed standard calculus.

\begin{theorem}[Normalization for intuitionistic natural deduction]
Every derivation in the calculus just fixed reduces to a normal derivation with
the same conclusion and no undischarged assumption that was not already
undischarged in the original derivation.  Equivalently, after identifying
alpha-equivalent formulas, the resulting set of undischarged assumption
formulas is a subset of the original set.
\end{theorem}

\paragraph{Proof sketch.}
Define maximal formulas and maximal segments.  Principal reductions remove an
introduction--elimination detour; permutation conversions move eliminations
above disjunction or existential eliminations.  A naive count of formula
occurrences is insufficient because a reduction can duplicate a
subderivation.  The standard proof uses a lexicographic measure on the rank
of maximal formulas and the lengths of maximal segments, or translates to a
typed term calculus and proves strong normalization there.  The full
calculus-specific proof is given in \cite{TroelstraSchwichtenberg2000}.

Normalization and cut elimination express related proof transformations, but
they are not one theorem.  Normalization concerns a reduction relation on
natural-deduction derivations; cut elimination concerns removal of a sequent
rule.

